%\ifodd\value{page}\hbox{}\newpage\fi
\chapter*{Śrī Lalitā Sahasranāma — Śloka 10}
\addcontentsline{toc}{chapter}{Śloka 10 - Nomi 22,23}

\begin{sloka}
{\sanskritfont
\begin{tabular}{c}
शुद्धविद्याङ्कुराकारद्विजपङ्क्तिद्वयोज्ज्वला ।\\
कर्पूरवीटिकामोदसमाकर्षिदिगन्तरा ॥ १०॥
\end{tabular}

\vspace{1.5em}

{\middlesize \iastfont
śuddha-vidyāṅkura-ākāra-dvija-paṅkti-dvaya-ujjvalā |\\
karpūra-vīṭikā-amoda-samākarṣi-dig-antarā ||10||}

\vspace{1em}

{\middlesize
\anvaya{%
शुद्धविद्याङ्कुराकारद्विजपङ्क्तिद्वयोज्ज्वला ।
कर्पूरवीटिकामोदसमाकर्षिदिगन्तरा ।
}
}

\middlesize \iastfont \itmean{%
«I cui due filari di denti, splendenti, hanno la forma di germogli della conoscenza pura;  
la fragranza della \textit{vīṭikā} al canforo trae a sé tutte le direzioni.»%
}

}
\end{sloka}

\wordmeaning{%
\wordline{शुद्ध-विद्या-अङ्कुर-आकार-द्विज-पङ्क्ति-द्वय-उज्ज्वला}
{śuddha-vidyā-aṅkura-ākāra-dvija-paṅkti-dvaya-ujjvalā}
{ splendente (\textit{ujjvalā}) per i due filari (\textit{paṅkti-dvaya}) di denti (\textit{dvija}) aventi la forma (\textit{ākāra}) di germogli (\textit{aṅkura}) della conoscenza pura (\textit{śuddha-vidyā}); }%
\wordline{कर्पूर-वीटिका-अमोद-समाकर्षि-दिक्-अन्तरा}
{karpūra-vīṭikā-amoda-\\samākarṣi-dig-antarā}
{ la fragranza (\textit{amoda}) della \textit{vīṭikā} al canforo (\textit{karpūra-vīṭikā}) trae a sé (\textit{samākarṣi}) gli spazi delle direzioni (\textit{dig-antarā}); }%
}

\textit{Śloka} 10 prosegue la descrizione della bocca di \textit{Lalitā}, ma spostando l’asse dal puro splendore sensibile al valore simbolico della \textit{vidyā}. I denti non sono soltanto bianchi e luminosi: essi sono detti “germogli della conoscenza pura”, indicando una parola che nasce da una sorgente non contaminata.

Nella tradizione di \textit{Śrīvidyā}, questi “germogli” evocano sia gli inizi della \textit{śuddha-vidyā} sia il mantra stesso, che, come un seme, si apre in suoni articolati attraverso la bocca della Dea. Si allude anche alla pienezza della trasmissione: i trentadue denti sono letti come simbolo della totalità delle iniziazioni e delle articolazioni del mantra, che dalla radice interiore giunge alla parola pronunciata.

Il secondo \textit{pāda} introduce l’elemento olfattivo. La fragranza della \textit{vīṭikā} al canforo non resta confinata al corpo, ma “trae a sé le direzioni”, cioè richiama le potenze che le governano, come se ogni punto dello spazio fosse desideroso di avvicinarsi alla sua presenza. Qui “spazi delle direzioni” rende \textit{dig-antarā} non come semplice orientamento astratto, ma come pervasività cosmica. La bellezza di \textit{Lalitā} non è statica né locale: essa irradia e convoca.

Questo \textit{śloka} contiene due \textit{nāma}.  
\textbf{\textit{Śuddha-vidyāṅkura-ākāra-dvija-paṅkti-dvaya-ujjvalā}} presenta i denti come manifestazione germinale della \textit{śuddha-vidyā}, luogo in cui la conoscenza pura prende forma e si prepara a farsi parola.  
\textbf{\textit{Karpūra-vīṭikā-amoda-samākarṣi-dig-antarā}} esprime la forza attrattiva della sua fragranza, che permea tutte le direzioni e convoca a sé l’intero orizzonte percettivo.
